% Construction : v(t) = (12 + t)e^{−0,6t} sur [0 ; 8], ordonnées divisées par 3.
% L'aire hachurée, de 0 à t, est la distance parcourue pendant le freinage ;
% elle croît vers 205/9 \(\approx\) 22,8 mètres, longueur minimale de la zone d'arrêt.
\begin{tikzpicture}[xscale=1.05, yscale=0.85]
  \draw[->] (0, 0) -- (8.6, 0);
  \draw[->] (0, 0) -- (0, 4.6) node[left] {$v$};
  \draw[thin] (0.300, 0) -- (0.300, 3.425);
  \draw[thin] (0.600, 0) -- (0.600, 2.930);
  \draw[thin] (0.900, 0) -- (0.900, 2.506);
  \draw[thin] (1.200, 0) -- (1.200, 2.142);
  \draw[thin] (1.500, 0) -- (1.500, 1.830);
  \draw[thin] (1.800, 0) -- (1.800, 1.562);
  \draw[thin] (2.100, 0) -- (2.100, 1.333);
  \draw[thin] (2.400, 0) -- (2.400, 1.137);
  \draw[thin] (2.700, 0) -- (2.700, 0.970);
  \draw[thin] (3.000, 0) -- (3.000, 0.826);
  \draw[thin] (3.300, 0) -- (3.300, 0.704);
  \draw[thin] (3.600, 0) -- (3.600, 0.600);
  \draw[thin] (3.900, 0) -- (3.900, 0.511);
  \draw[thin] (4.200, 0) -- (4.200, 0.434);
  \draw[thick] plot coordinates {(0.000, 4.000) (0.167, 3.670) (0.333, 3.366) (0.500, 3.087) (0.667, 2.830) (0.833, 2.595) (1.000, 2.378) (1.167, 2.179) (1.333, 1.997) (1.500, 1.830) (1.667, 1.676) (1.833, 1.535) (2.000, 1.406) (2.167, 1.287) (2.333, 1.178) (2.500, 1.078) (2.667, 0.987) (2.833, 0.903) (3.000, 0.826) (3.167, 0.756) (3.333, 0.692) (3.500, 0.633) (3.667, 0.579) (3.833, 0.529) (4.000, 0.484) (4.167, 0.442) (4.333, 0.404) (4.500, 0.370) (4.667, 0.338) (4.833, 0.309) (5.000, 0.282) (5.167, 0.258) (5.333, 0.236) (5.500, 0.215) (5.667, 0.197) (5.833, 0.180) (6.000, 0.164) (6.167, 0.150) (6.333, 0.137) (6.500, 0.125) (6.667, 0.114) (6.833, 0.104) (7.000, 0.095) (7.167, 0.087) (7.333, 0.079) (7.500, 0.072) (7.667, 0.066) (7.833, 0.060) (8.000, 0.055)};
  \draw[dotted] (4.5, 0) -- (4.5, {(12 + 4.5) * exp(-2.7) / 3});
  \node[below] at (4.5, -0.05) {$t$};
  \node[right] at (5.4, 1.1) {$\mathcal{C}_v$};
\end{tikzpicture}
